% !TEX root = ../main.tex

\chapter{引言}
\label{ch:intro}

\section{研究背景与意义}
时空转录组学（Spatial Transcriptomics, ST）作为近年来生命科学领域最具突破性的技术之一，成功将高通量单细胞测序（scRNA-seq）的基因表达定量能力与细胞在组织中的空间位置信息融为一体\cite{Asp2020}。通过诸如 Stereo-seq、Visium 等高分辨率时空测序技术，研究人员能够在不损失空间结构的条件下，观测多细胞生物器官发生和发育过程中成千上万个基因的动态表达模式\cite{Chen2022}。然而，这类测序通常是“破坏性”的，即每个样本在测序时均已被固定，无法对同一个活体细胞进行跨时间点的连续追踪。因此，如何在多个离散时间采样点（Time-series）的空间转录组数据中，精准重建细胞的迁移路线和谱系分化轨迹，成为了发育生物学和计算生物学亟待解决的核心科学问题。

\section{现有谱系追踪算法及局限性分析}
近年来，计算生物学家尝试将数学中的最优传输（Optimal Transport, OT）理论引入单细胞动力学建模中\cite{Villani2009}。其中最具代表性的工作包括基于静态表达数据重建细胞发育轨迹的 WOT（Waddington Optimal Transport）\cite{Schiebinger2019}，以及近年来将最优传输扩展到多模态和空间数据的 Moscot（Multi-Omic Single-Cell Optimal Transport）\cite{Klein2024}。尽管这些方法在解析单细胞表达谱流形方面取得了一定进展，但当面对时空转录组数据中复杂的二维或三维空间拓扑结构时，它们普遍存在以下两个关键痛点：

\begin{enumerate}
    \item \textbf{欧氏空间瞬间移动（Euclidean Teleportation）}：
    传统的空间最优传输模型在衡量细胞空间迁移代价时，默认采用欧氏空间距离（Euclidean Distance）作为物理邻近性的度量。这一假设在组织拓扑结构简单的区域是适用的，但在包含复杂几何屏障（例如小鼠胚胎发育过程中的脑室、发育空腔或体外培养形成的物理空隙）时，欧氏距离会发生失效。由于欧氏距离走的是两点间的“直线”，它会忽视组织实际的几何轮廓，导致预测出的细胞迁移路径直接“穿透”发育腔洞等物理真空区域。这种物理上不可能发生的“瞬间移动”伪影严重损害了谱系重建的物理真实性。
    
    \item \textbf{时间方向对称与缺乏方向性约束（Directional Symmetry）}：
    标准的无向最优传输模型在计算最优耦合矩阵时，其传输代价（Cost）关于传输方向是对称的，即从细胞 $i$ 传输到细胞 $j$ 的代价与反向传输是等同的。虽然部分方法引入了伪时间（Pseudotime）约束，但仍无法精细捕捉局部的细胞分化动力学箭头。实际上，发育过程具有极强的方向性。单细胞 RNA 速度（RNA Velocity）作为转录动力学的微观表现形式，能够指示每个细胞在转录流形上的瞬间运动方向\cite{LaManno2018, Bergen2020}。常规的最优传输模型因无法将这种局部非对称的动态速度矢量融入传输代价中，导致在复杂的分化分叉点极易产生预测失真，无法准确捕捉细胞谱系的时间方向性。
\end{enumerate}

\section{本文主要研究内容与创新点}
针对上述局限性，本文提出了一个物理和几何双重约束的时空单细胞谱系追踪框架——\textbf{SpaLineage-OT}。其核心目标是克服“欧氏空间瞬间移动”与“方向对称性”瓶颈，建立满足组织拓扑约束与生物时间箭头的连续时空分化动力学模型。

本文的创新点与核心贡献主要包括：
\begin{enumerate}
  \item \textbf{黎曼测地线几何约束}：
  我们抛弃了传统的欧氏空间假设，利用 Delaunay 三角剖分算法将空间细胞坐标构建为空间流形网格图，并引入基于核密度估计（KDE）的形态阻抗函数，将空间流形映射为非均匀的黎曼流形。通过 Dijkstra 算法求解该流形上的全源最短路径，获取黎曼测地线距离（Riemannian Geodesic Distance）作为最优传输的空间代价。这强制了细胞最优传输耦合必须紧密贴合组织流形，彻底消除了跨越物理屏障的“瞬间移动”伪影。
  
  \item \textbf{非对称速度引导传输代价（AS-FGW）}：
  我们将单细胞 RNA 速度向量投射至物理空间，计算细胞间几何位移向量与速度矢量的夹角余弦相似度，并将其作为非对称方向惩罚项引入融合格罗莫夫-瓦瑟斯坦（FGW）最优传输框架中。由此，我们构建了首个将 RNA 速度与空间拓扑结合的非对称有向最优传输耦合矩阵。
  
  \item \textbf{薛定谔桥流匹配与连续动力学拟合}：
  利用最优传输耦合作为桥接，我们将离散时间采样点数据转化为生成式流匹配问题。通过构建薛定谔桥流匹配（Schrödinger Bridge Flow Matching, SBFM）框架\cite{Lipman2023, Liu2023}，并使用多层感知机（DriftMLP）拟合连续时间的动力学漂移场（Drift Field）。为了约束流场使其满足物理守恒定律，我们进一步将局部 RNA 速度作为物理先验（Physics-Prior Loss）融入神经网络训练中。
  
  \item \textbf{均值漂移势能引导与常微分方程积分}：
  在利用四阶龙格-库塔（RK4）方法求解 Neural ODE 进行细胞轨迹推断时，由于复杂的边缘区域极易发生数值下溢，我们基于目标时间点的细胞空间分布密度构建了均值漂移势能场（Mean Shift Potential Field）。其负梯度提供了一个指向组织内部的高维引力约束，纠正了边缘细胞的发散漂移，确保了演化路径的数值稳定度与物理合理性。
\end{enumerate}

\section{论文结构安排}
本学位论文的结构安排如下：
第一章（本章）介绍了时空转录组谱系追踪的研究背景、现有方法瓶颈以及本文的方法概述与创新点。
第二章将深入阐述 SpaLineage-OT 模型的数学原理，详细推导黎曼测地线约束、非对称速度引导代价、SBFM 流匹配框架以及均值漂移势能的公式。
第三章将介绍基于真实小鼠胚胎器官发育时空数据集（MOSTA）的实验设计、数据预处理以及计算流程。
第四章将展示实验结果，通过多维度的定量指标（能量距离、屏障穿越率等）与可视化图表（Figure 1-8 配合架构流程图 Figure 9）与 Moscot、WOT 等基线方法进行对比。
第五章将对模型的物理局限性进行讨论，并展望未来的计算加速与多模态扩展。
最后，我们将总结全文。
